% TOP_PROBLEM: {"id": 30006974, "title": "Irrationality of e+pi", "category_id": 1, "category": {"id": 1, "name": "number_theory", "display_name": "Number Theory", "description": "Properties of integers, prime numbers, Diophantine equations.", "slug": "number-theory", "order_index": 1, "created_at": "2026-07-31T15:26:25.670Z"}, "status": "open"}
% FIRST_POSTED: null
% ATTEMPT_STATUS: unresolved
% Consolidated research record by GPT-6 Astra Pro, 27 September 2026.
% Mathematical corrections and supplementary checks are explicitly marked below.
% Compile with LuaLaTeX or XeLaTeX, twice. No external bibliography file is needed.
\documentclass[11pt]{article}
\usepackage[margin=25mm]{geometry}
\usepackage{fontspec}
\setmainfont{Latin Modern Roman}
\usepackage{amsmath,amssymb,mathtools}
\usepackage{unicode-math}
\setmathfont{Latin Modern Math}
\usepackage{xurl,hyperref}
\hypersetup{colorlinks=true,urlcolor=blue,
  pdftitle={Irrationality of e+pi: consolidated research attempt},
  pdfauthor={GPT-6 Astra Pro},
  pdfsubject={Conditional criteria, factorial identities, transcendence routes, and unresolved gaps}}
\setcounter{secnumdepth}{0}
\setlength{\parindent}{0pt}
\setlength{\parskip}{5pt}
\setlength{\emergencystretch}{3em}
\newcommand{\Q}{\mathbb Q}
\newcommand{\Z}{\mathbb Z}
\newcommand{\R}{\mathbb R}
\newcommand{\C}{\mathbb C}
\newcommand{\Qbar}{\overline{\mathbb Q}}
\DeclareMathOperator{\trdeg}{trdeg}
\DeclareMathOperator{\Log}{Log}
\DeclareMathOperator{\ord}{ord}
\begin{document}
% problemId: 30006974
% Canonical problem ID: 30006974
\section*{\texorpdfstring{Irrationality of e+pi}{Irrationality of e+pi}}\label{30006974}
\label{problem:pi-plus-e}

\subsection{Definitions and mathematical statement}
Let
\[
 e=\exp(1)=\sum_{k=0}^{\infty}\frac1{k!},\qquad
 \pi=4\arctan(1)=4\int_0^1\frac{dt}{1+t^2}.
\]
Here $0!=1$, and $n!=1\cdot2\cdots n$ for $n\ge1$. A real number is
\emph{rational} if it belongs to
$\Q=\{a/b:a,b\in\Z,\ b>0\}$, and is \emph{irrational} otherwise.
The exact question is to prove or disprove
\begin{equation}\label{eq:target}
 \boxed{e+\pi\notin\Q.}
\end{equation}
Equivalently, does
\[
 \forall a\in\Z\ \forall b\in\Z_{>0},\qquad b(e+\pi)\ne a
\]
hold? A disproof must supply a rational value with an exact proof of the
equality; a finite decimal approximation would not suffice.

A complex number is \emph{algebraic} if it is a zero of a nonzero polynomial
in $\Q[T]$; $\Qbar$ denotes the field of all algebraic numbers. A number
outside $\Qbar$ is \emph{transcendental}. Two numbers $x,y$ are
\emph{algebraically independent over $\Q$} if no nonzero
$P\in\Q[X,Y]$ satisfies $P(x,y)=0$.
The transcendence of $e$ and of $\pi$ separately is established background,
not a solution of \eqref{eq:target}; see
\href{https://webusers.imj-prg.fr/~michel.waldschmidt/articles/pdf/ColloquiumDeGiorgiSchanuelConjecture2014.pdf}{[E5]}.
The stronger assertions $e+\pi\notin\Qbar$, rational linear independence
of $1,e,\pi$, and algebraic independence of $e,\pi$ are kept distinct
from the selected target throughout this record.

\par\smallskip\textit{Formulation and status references:}
\href{https://arxiv.org/html/2311.00546v2}{[S1]}, Section 2.1;
\href{https://arxiv.org/abs/2606.17303}{[S2]}, abstract.
These are dated source checkpoints, not substitutes for a proof.

\subsection{Short English statement}
Can the sum of Euler's number and the circle constant be a fraction of two
integers? Establish that no such fraction exists, or exhibit one with a
rigorous proof of equality.

\subsection{Research attempt}

\paragraph{Scope, provenance, and source checkpoint.}
This is a consolidated mathematical export of the two earlier research
responses in this conversation and the additional development supplied by
the user. It preserves their arguments, proved implications, unsuccessful
routes, and remaining gaps. Repeated explanations have been combined;
mathematical corrections are identified explicitly rather than silently
incorporated. The uploaded Birch--Swinnerton-Dyer document supplies the
format only, not mathematical evidence about $e+\pi$.

The first response established a denominator-matching criterion. The second
proved forbidden intervals for good approximations to $\pi$ under the
rationality hypothesis, and deduced superexponential growth of the
continued-fraction denominators. The supplied continuation adds the
exponential, factorial, Schanuel, and $E/G$-function routes, together with
an examination of a claimed proof. The proofs below are ordinary written
arguments; no Lean formalization or independent peer review is asserted.
No claim of priority or literature-wide novelty is made.

Kowalski's 2025 revision explicitly records the irrationality question as
unresolved \href{https://arxiv.org/html/2311.00546v2}{[S1]}.
The June 2026 preprint cited in the first response also describes it as
open \href{https://arxiv.org/abs/2606.17303}{[S2]}.
The source checks for this export were performed on 27 September 2026.
The criticism of \href{https://arxiv.org/html/2007.15000v2}{[D1]} below
concerns the explicitly cited version~2; it is not presented as an audit
of every subsequent version.

\paragraph{Verdict after a serious attempt.}
I tried to derive a contradiction from $e+\pi\in\Q$ using compatible
rational approximations, exponential identities, factorial tails,
transcendence theorems, and the arithmetic theory of $E$- and $G$-functions.
I did not obtain an unconditional proof or disproof. What has been
established here consists of exact consequences of rationality and fully
proved sufficient conditions for irrationality. None of the missing
conditions is assumed to have been established for the actual number
$\pi$.

\paragraph{1. The common rationality hypothesis.}
Whenever a contradiction argument is being attempted, suppose
\begin{equation}\label{eq:rationality}
 s=e+\pi=\frac ab,\qquad a,b\in\Z,\quad b>0.
\end{equation}
Thus $\pi=s-e$. The symbols $a,b$ are fixed throughout any one such
argument. It is not necessary to assume that $a/b$ is reduced. Bounds that
depend on $b$ are permissible, but taking $b$ to be bounded without proof
would not be permissible.

\paragraph{2. Continued-fraction notation and the special approximations to $e$.}
For an irrational real $x=[c_0;c_1,c_2,\ldots]$, its regular
continued-fraction convergents $h_j/k_j$ are reduced and satisfy
\[
 h_j=c_jh_{j-1}+h_{j-2},\qquad
 k_j=c_jk_{j-1}+k_{j-2},
\]
with $h_{-2}=0$, $h_{-1}=1$, $k_{-2}=1$, and $k_{-1}=0$.
In particular, $k_j>0$ for $j\ge0$.
If $\xi_{j+1}=[c_{j+1};c_{j+2},\ldots]$, the determinant identity and
\[
 x=\frac{h_j\xi_{j+1}+h_{j-1}}{k_j\xi_{j+1}+k_{j-1}}
\]
give
\begin{equation}\label{eq:cf-error}
 0<\left|x-\frac{h_j}{k_j}\right|
 =\frac1{k_j(k_j\xi_{j+1}+k_{j-1})}
 <\frac1{c_{j+1}k_j^2},\qquad
 \left|x-\frac{h_j}{k_j}\right|<\frac1{k_jk_{j+1}}.
\end{equation}
These are the standard tools used here; see
\href{https://arxiv.org/html/math/0601660v3}{[E1]} and
\href{https://wstein.org/edu/2007/spring/ent/ent-html/node65.html}{[E2]}.

Euler's explicit expansion is
\begin{equation}\label{eq:e-cf}
 e=[2;1,2,1,\ 1,4,1,\ 1,6,1,\ldots].
\end{equation}
Cohn gives a proof, rather than a numerical extrapolation, of this
identity \href{https://arxiv.org/html/math/0601660v3}{[E1]}.
Let $p_n/q_n$ be the convergent that stops immediately before the partial
quotient $2n$. With the initial digit indexed by zero, this is convergent
$3n-2$. Thus
\[
 \frac{p_1}{q_1}=\frac31,\qquad
 \frac{p_2}{q_2}=\frac{19}{7},\qquad
 \frac{p_3}{q_3}=\frac{193}{71},\qquad
 \frac{p_4}{q_4}=\frac{2721}{1001}.
\]
The next digit in \eqref{eq:e-cf} and \eqref{eq:cf-error} imply
\begin{equation}\label{eq:e-small}
 0<\left|e-\frac{p_n}{q_n}\right|<\frac1{2nq_n^2},
 \qquad
 \boxed{|q_ne-p_n|<\frac1{2nq_n}.}
\end{equation}
The extra factor $n$ is the resource that the first argument attempts to
combine with good approximations to $\pi$.

\paragraph{3. First proved criterion: compatible approximation denominators.}
Consider the following specific condition.

\textbf{Condition $(*)$.} For infinitely many positive integers $n$, there
exists a reduced fraction $u_n/v_n$, with $v_n>0$, such that
\begin{equation}\label{eq:matching}
 n^{1/3}q_n\le v_n\le n^{2/3}q_n,
 \qquad
 \left|\pi-\frac{u_n}{v_n}\right|<\frac1{v_n^2}.
\end{equation}

\textbf{Proposition 1.} Condition $(*)$ implies $e+\pi\notin\Q$.

\emph{Proof.} Assume \eqref{eq:rationality}, and, for an $n$ supplied by
$(*)$, put
\begin{equation}\label{eq:matching-form}
 L_n=q_nv_n(e+\pi)-(p_nv_n+u_nq_n).
\end{equation}
Then
\[
 bL_n=aq_nv_n-bp_nv_n-bu_nq_n\in\Z.
\]
On the other hand,
\[
 L_n=v_n(q_ne-p_n)+q_n(v_n\pi-u_n).
\]
Consequently,
\begin{align}
 |L_n|
 &\le v_n|q_ne-p_n|+q_n|v_n\pi-u_n|\notag\\
 &<\frac{v_n}{2nq_n}+\frac{q_n}{v_n}
 \le\frac1{2n^{1/3}}+\frac1{n^{1/3}}
 =\frac3{2n^{1/3}}.\label{eq:matching-small}
\end{align}
Along the infinite set of admissible $n$, this tends to zero. Hence for
sufficiently large such $n$, $|bL_n|<1$ and the integer $bL_n$ must be zero.

It is essential to rule out exact cancellation. If $L_n=0$, then
\[
 \frac{u_n}{v_n}=\frac ab-\frac{p_n}{q_n}
 =\frac{aq_n-bp_n}{bq_n}.
\]
Since $u_n/v_n$ is reduced, $v_n$ divides $bq_n$, so $v_n\le bq_n$.
But \eqref{eq:matching} gives $v_n\ge n^{1/3}q_n>bq_n$ whenever $n>b^3$.
This contradiction proves the proposition.\hfill$\square$

\paragraph{4. Why the first criterion is not yet a solution.}
The convergents of $\pi$ supply infinitely many reduced fractions $u/v$
with $|\pi-u/v|<1/v^2$, by \eqref{eq:cf-error}. That fact alone does not
place any of their denominators in the prescribed intervals
\[
 [n^{1/3}q_n,\ n^{2/3}q_n]
\]
for infinitely many $n$. The missing assertion is precisely this
compatibility of scales. The ratio $v_n/q_n$ must tend to infinity to
exclude rational cancellation, but must stay sufficiently small compared
with $n$ for the first term in \eqref{eq:matching-small} to tend to zero.

A decisive logical stress test is $x=6-e$. This number is irrational
(in fact transcendental), while $e+x=6$. The proof of Proposition~1 works
unchanged with $x$ in place of $\pi$. Therefore $(*)$ must fail for $x$.
In particular, $(*)$ cannot be inferred merely from the irrationality of
the second summand, or from the existence of infinitely many good
approximations to each summand separately.

\paragraph{5. The selected denominator recurrence.}
For later use, extend the notation by $q_0=1$. The full convergent
recurrences yield
\begin{equation}\label{eq:q-recurrence}
 \boxed{q_{n+1}=(4n+2)q_n+q_{n-1},\qquad q_0=q_1=1.}
\end{equation}
Here is the three-step calculation. If $r_n$ is the full convergent
denominator immediately before $q_n$, the next three digits are $2n,1,1$.
Their denominators are
\[
 2nq_n+r_n,\qquad (2n+1)q_n+r_n,\qquad (4n+1)q_n+2r_n.
\]
The two preceding digits equal to $1$ give $q_n=2r_n-q_{n-1}$ for $n\ge2$;
for $n=1$ this identity holds with $r_1=q_0=1$. Substitution proves
\eqref{eq:q-recurrence}.

It follows that
\begin{equation}\label{eq:q-bounds}
 q_n\ge n!,\qquad q_{n+1}\le(4n+3)q_n\quad(n\ge1).
\end{equation}
Indeed, the denominators are nondecreasing, so $q_{n-1}\le q_n$ proves the
upper bound. For $n\ge2$, the recurrence gives
$q_n\ge(4n-2)q_{n-1}\ge n q_{n-1}$; induction from $q_1=1$ proves
the lower bound.

\paragraph{6. Second proved result: forbidden denominator intervals.}
\textbf{Proposition 2.} Under \eqref{eq:rationality}, for every integer
$n\ge2b^2$ there is no reduced fraction $u/v$, $v>0$, satisfying both
\begin{equation}\label{eq:forbidden-range}
 2bq_n\le v\le\frac{nq_n}{b}
\end{equation}
and
\begin{equation}\label{eq:good-pi}
 \left|\pi-\frac uv\right|<\frac1{v^2}.
\end{equation}

\emph{Proof.} Define
\[
 r_n^{*}=\frac ab-\frac{p_n}{q_n}
 =\frac{aq_n-bp_n}{bq_n}.
\]
The superscript distinguishes this rational number from the denominator
$r_n$ used only in the previous calculation. By \eqref{eq:e-small},
\begin{equation}\label{eq:shifted-error}
 |\pi-r_n^{*}|<\frac1{2nq_n^2}.
\end{equation}
The reduced denominator of $r_n^{*}$ is at most $bq_n$, so
\eqref{eq:forbidden-range} implies $u/v\ne r_n^{*}$. Therefore
\[
 \left|r_n^{*}-\frac uv\right|
 =\frac{|(aq_n-bp_n)v-bq_nu|}{bq_nv}
 \ge\frac1{bq_nv}.
\]
Combining this with the triangle inequality, \eqref{eq:shifted-error},
and \eqref{eq:good-pi} gives
\[
 1<\frac{bv}{2nq_n}+\frac{bq_n}{v}.
\]
The upper and lower bounds in \eqref{eq:forbidden-range}, respectively,
make the two terms on the right at most $1/2$. This is impossible.
\hfill$\square$

Thus rationality would exclude every sufficiently good approximation to
$\pi$ whose denominator belongs to
\begin{equation}\label{eq:forbidden-interval}
 I_n=\left[2bq_n,\frac{nq_n}{b}\right],\qquad n\ge2b^2.
\end{equation}
The ratio of its endpoints is $n/(2b^2)$ and tends to infinity.

\paragraph{7. Third proved result: superexponential growth would follow.}
Write $P_j/Q_j$ for the successive regular continued-fraction convergents
of $\pi$, indexed by $j\ge0$. Changing finitely many starting indices has
no effect on the following limit.

\textbf{Proposition 3.} If $e+\pi\in\Q$, then
\begin{equation}\label{eq:superexponential}
 \boxed{\lim_{j\to\infty}\frac{\log Q_j}{j}=+\infty.}
\end{equation}
Equivalently, for every fixed $C>1$, eventually $Q_j>C^j$.

\emph{Proof.} Every $P_j/Q_j$ satisfies \eqref{eq:good-pi}, so its
denominator avoids all the intervals \eqref{eq:forbidden-interval}.
These intervals occur in increasing order, and
$nq_n/b<2bq_{n+1}$ by \eqref{eq:q-recurrence}. Consequently all sufficiently
large $Q_j$ lie in one of the gaps
\begin{equation}\label{eq:gaps}
 J_n=\left(\frac{nq_n}{b},\ 2bq_{n+1}\right),\qquad n\ge2b^2.
\end{equation}
Their endpoint ratios are uniformly bounded:
\begin{equation}\label{eq:gap-ratio}
 \frac{2bq_{n+1}}{nq_n/b}
 \le2b^2\frac{4n+3}{n}\le14b^2.
\end{equation}
Also the convergent recurrence for $\pi$ implies
\begin{equation}\label{eq:two-step}
 Q_{j+2}\ge Q_{j+1}+Q_j\ge2Q_j.
\end{equation}
Thus an interval with endpoint ratio at most $14b^2$ contains at most
\[
 C_b=2\left\lceil\log_2(14b^2)\right\rceil+3
\]
convergent denominators. To see this directly, among more than $C_b$
consecutive denominators in such an interval, the first and one of the
last denominators would be separated by enough two-step doublings in
\eqref{eq:two-step} to exceed that ratio.

The finitely many initial denominators can be absorbed in a constant
$j_0$. Counting at most $C_b$ denominators per gap then shows that, whenever
$Q_j\in J_n$,
\begin{equation}\label{eq:index-bound}
 j\le j_0+C_bn.
\end{equation}
Meanwhile \eqref{eq:gaps} and \eqref{eq:q-bounds} give
\[
 Q_j>\frac{nq_n}{b}\ge\frac{n\,n!}{b}.
\]
For sufficiently large $n$ the logarithm of this last bound is positive,
and hence
\begin{equation}\label{eq:growth-ratio}
 \frac{\log Q_j}{j}
 \ge\frac{\log(n\,n!/b)}{j_0+C_bn}\longrightarrow+\infty.
\end{equation}
For example, at least $\lfloor n/2\rfloor$ factors in $n!$ are at least
$n/2$, so the numerator has growth at least of order $n\log n$.
Finally $n\to\infty$ as $j\to\infty$, because each of the first finitely
many gaps is bounded and contains finitely many denominators. This proves
the asserted limit.\hfill$\square$

\textbf{Corollary.} An unconditional proof that some fixed $C>1$ satisfies
\begin{equation}\label{eq:missing-growth}
 Q_j\le C^j\quad\text{for infinitely many }j
\end{equation}
would prove $e+\pi\notin\Q$.
This is a sufficient condition; the converse is not asserted.

\paragraph{8. The remaining growth gap and irrationality measures.}
No proof of \eqref{eq:missing-growth} for $\pi$ has been supplied in this
attempt. Superexponential growth itself is not a contradiction: irrational
numbers may have continued fractions with arbitrarily large partial
quotients. The model $6-e$ from paragraph~4 is also compatible with
Proposition~3, whose proof uses no special property of $\pi$ beyond its
irrationality and the assumed relation to $e$.

For an irrational $x$, define its irrationality exponent by
\[
 \mu(x)=\sup\left\{\mu>0:
 0<\left|x-\frac pq\right|<q^{-\mu}
 \text{ for infinitely many reduced }p/q,\ q>0\right\}.
\]
Zeilberger and Zudilin proved the bound
\begin{equation}\label{eq:pi-measure}
 \mu(\pi)\le7.103205334137\ldots
\end{equation}
\href{https://arxiv.org/abs/1912.06345}{[E7]}.
Only this published bound is used; no claim that it is the latest possible
bound is required here.

For any fixed $M$ larger than the bound in \eqref{eq:pi-measure}, all
sufficiently large convergents satisfy
\[
 Q_j^{-M}\le\left|\pi-\frac{P_j}{Q_j}\right|
 <\frac1{Q_jQ_{j+1}}.
\]
Thus one obtains only
\begin{equation}\label{eq:power-growth}
 Q_{j+1}<Q_j^{M-1}.
\end{equation}
This controls growth from one denominator to the next as a power of that
denominator, not as an exponential function of its index. It is compatible,
for instance, with logarithms growing like $j\log j$. Therefore this
irrationality-measure argument does not contradict
\eqref{eq:superexponential} and does not establish
\eqref{eq:missing-growth}.

\paragraph{9. The exponential attack.}
Under \eqref{eq:rationality}, Euler's identity gives
\[
 -1=\exp(i\pi)=\exp\bigl(i(s-e)\bigr)=\exp(is)\exp(-ie),
\]
so
\begin{equation}\label{eq:exponential}
 \boxed{\exp(ie)=-\exp(is).}
\end{equation}
Because $s=e+\pi>0$ and is assumed rational, $1$ and $is$ are algebraic and
linearly independent over $\Q$. The Lindemann--Weierstrass theorem states
that exponentials of $\Q$-linearly independent algebraic numbers are
algebraically independent over $\Qbar$; see
\href{https://arxiv.org/html/2406.01296v1}{[E3]}, Theorem~3.
It follows that $e=\exp(1)$ and $\exp(is)$ are algebraically independent.
By \eqref{eq:exponential}, the same would then be true of $e$ and $\exp(ie)$.

For clarity, the algebraic independence conclusion also follows by
expanding a hypothetical polynomial relation in $e$ and $\exp(is)$.
Its monomials are exponentials $\exp(m+nis)$ with distinct algebraic
exponents, because $1$ and $is$ are $\Q$-linearly independent. The
linear-independence form of Lindemann--Weierstrass excludes a nontrivial
algebraic linear relation among them.

This is a valid consequence of rationality, but no contradiction has been
derived. The exponent $ie$ is transcendental, and the quoted theorem does
not directly apply to that exponent. In particular, one may not extend its
hypotheses from algebraic exponents to arbitrary transcendental ones.

Taking logarithms of \eqref{eq:exponential} only gives
\[
 e-s\equiv\pi\pmod{2\pi}.
\]
The original relation $e-s=-\pi$ already satisfies this congruence. The
congruence loses branch information rather than producing a new
arithmetic obstruction. With the specified logarithm value
$\Log(-1)=i\pi$, one can write
\begin{equation}\label{eq:mixed-exp-log}
 e+\pi=\exp(1)-i\Log(-1).
\end{equation}
Here $\Log(-1)=i\pi$ denotes a chosen value; no analytic branch across a
branch cut is being assumed. Excluding an algebraic value of this
particular exponential--logarithm sum would suffice, but no such theorem
is proved by the displayed identities.

\paragraph{10. The factorial-tail attack.}
Define the integers and positive remainders
\begin{equation}\label{eq:factorial-setup}
 A_n=\sum_{k=0}^n\frac{n!}{k!}\in\Z,\qquad
 R_n=n!e-A_n.
\end{equation}
For every $n\ge1$,
\[
 R_n=\frac1{n+1}+\frac1{(n+1)(n+2)}
       +\frac1{(n+1)(n+2)(n+3)}+\cdots.
\]
The first term gives a strict lower bound, and comparison with a geometric
series gives the strict upper bound
\begin{equation}\label{eq:factorial-tail}
 \boxed{\frac1{n+1}<R_n<\sum_{j=1}^{\infty}\frac1{(n+1)^j}=\frac1n.}
\end{equation}
If $n\ge b$, then $n!s\in\Z$, so
\begin{equation}\label{eq:K-definition}
 K_n=n!s-A_n\in\Z,\qquad n!\pi=K_n-R_n.
\end{equation}
For a real number $x$, let $\{x\}=x-\lfloor x\rfloor$ and
$\|x\|=\min_{m\in\Z}|x-m|$. For $n\ge\max\{b,3\}$,
$0<R_n<1/2$, whence
\begin{equation}\label{eq:fingerprint}
 \boxed{\{n!\pi\}=1-R_n,\qquad \|n!\pi\|=R_n.}
\end{equation}
In particular, $K_n=\lceil n!\pi\rceil$, and division of
\eqref{eq:factorial-tail} by $n!$ gives
\begin{equation}\label{eq:factorial-approx}
 \boxed{\frac1{(n+1)!}<\frac{K_n}{n!}-\pi<\frac1{n\,n!}.}
\end{equation}
The common threshold $\max\{b,3\}$ makes explicit the two separate
conditions appearing in the supplied development.

\paragraph{11. Why the factorial approximations are not strong enough.}
With the displayed denominator $q=n!$, the error in
\eqref{eq:factorial-approx} is of order $1/(nq)$. Since
$\log n/\log(n!)\to0$, this has scale
\[
 q^{-1-o(1)},
\]
not $q^{-2}$ or a still stronger approximation exponent. Moreover $K_n/n!$
need not be reduced: nothing in this estimate controls its reduced
denominator. It therefore supplies no contradiction to a bound on the
irrationality exponent of $\pi$. Possible cancellation in a displayed
denominator cannot simply be assumed to make the approximations good
enough.

The integrality of $A_n$ follows term by term. The identities
\[
 A_{n+1}=(n+1)A_n+1,\qquad
 R_{n+1}=(n+1)R_n-1
\]
give
\begin{equation}\label{eq:K-recurrence}
 \boxed{K_{n+1}=(n+1)K_n-1.}
\end{equation}
This is again rigid, but not inconsistent.

\textit{Explicit clarification of the recurrence reduction.}
For any integer sequence satisfying \eqref{eq:K-recurrence} from $n=N$
onward,
\[
 \frac{K_{n+1}}{(n+1)!}=\frac{K_n}{n!}-\frac1{(n+1)!}.
\]
Telescoping and letting $n\to\infty$ yields
\begin{equation}\label{eq:recurrence-limit}
 \lim_{n\to\infty}\frac{K_n}{n!}
 =\frac{K_N}{N!}-\sum_{k=N+1}^{\infty}\frac1{k!}
 =\frac{K_N+A_N}{N!}-e.
\end{equation}
The limit is precisely a rational number minus $e$. Thus proving the
recurrence for $K_n=\lceil n!\pi\rceil$ eventually would already imply
$e+\pi\in\Q$, since $K_n/n!\to\pi$. Conversely rationality gives that
recurrence and that ceiling identity. This is a reformulation, not an
independent contradiction. The related factorial-arithmetic theme also
appears in \href{https://arxiv.org/abs/2606.17303}{[S2]}; the calculation
above is written out here.

\paragraph{12. An exact factorial-digit fingerprint.}
A factorial-base expansion of a real number is an integer part followed by
\[
 \sum_{m=2}^{\infty}\frac{d_m}{m!},\qquad
 d_m\in\{0,1,\ldots,m-1\}.
\]
Use the convention excluding tails eventually consisting entirely of the
maximal digits $m-1$; this removes the usual ambiguity at terminating
expansions. Every rational has a terminating expansion in this convention:
its denominator divides $N!$ for some $N$, and repeated integer division
in the radices $2,3,\ldots,N$ supplies the finite digits.

Since
\[
 e=2+\sum_{m=2}^{\infty}\frac1{m!},
\]
a rational value of $s=e+\pi$ would force an eventual tail with
$d_m=m-2$ for $\pi$. The borrowing identity is
\begin{equation}\label{eq:borrowing}
 -\sum_{m=N+1}^{\infty}\frac1{m!}
 =-\frac1{N!}+\sum_{m=N+1}^{\infty}\frac{m-2}{m!}.
\end{equation}
Indeed,
\[
 \sum_{m=N+1}^{\infty}\frac{m-1}{m!}
 =\sum_{m=N+1}^{\infty}
   \left(\frac1{(m-1)!}-\frac1{m!}\right)=\frac1{N!},
\]
which proves \eqref{eq:borrowing} by subtracting the factorial tail.
If $N!s$ is integral, then
\begin{equation}\label{eq:digit-tail}
 \pi=\frac{K_N-1}{N!}
       +\sum_{m=N+1}^{\infty}\frac{m-2}{m!}.
\end{equation}
The tail is positive and less than $1/N!$ for $N\ge2$. Consequently this
really specifies a finite prefix followed by admissible digits $m-2$;
it is not merely a formal rearrangement with invalid digits.

\textit{Explicit clarification of the converse.}
If $\pi$ has such a tail after a finite factorial prefix $B/N!$, then
\[
 \pi+e=\frac B{N!}+\sum_{m=0}^{N}\frac1{m!}+\frac1{N!}\in\Q.
\]
Thus the eventual digit pattern is equivalent to the rationality
hypothesis, under the stated expansion convention. No theorem about the
factorial digits of $\pi$ excluding this tail has been supplied here.
Every $r-e$, $r\in\Q$, genuinely has such a tail and is transcendental,
since otherwise $e=r-(r-e)$ would be algebraic. The transcendence of
$\pi$ by itself therefore does not rule out the fingerprint.

\paragraph{13. A complete conditional proof from Schanuel's conjecture.}
Schanuel's conjecture asserts that if $z_1,\ldots,z_m\in\C$ are linearly
independent over $\Q$, then
\begin{equation}\label{eq:schanuel}
 \trdeg_{\Q}\Q(z_1,\ldots,z_m,e^{z_1},\ldots,e^{z_m})\ge m.
\end{equation}
Here transcendence degree is the maximal cardinality of an algebraically
independent subset of the field over $\Q$. The statement and its standard
application to $e,\pi$ are discussed in
\href{https://arxiv.org/html/2403.09304v1}{[E4]}, Section~2.1, and
\href{https://webusers.imj-prg.fr/~michel.waldschmidt/articles/pdf/ColloquiumDeGiorgiSchanuelConjecture2014.pdf}{[E5]}.

\textbf{Proposition 4.} Schanuel's conjecture implies that $e+\pi$ is
transcendental.

\emph{Proof.} Choose $z_1=1$ and $z_2=i\pi$. They are $\Q$-linearly
independent: a rational relation has zero real part and zero imaginary
part only when both coefficients vanish. Since $e^{i\pi}=-1$,
\eqref{eq:schanuel} would give
\[
 \trdeg_{\Q}\Q(e,i\pi)\ge2.
\]
Adjoining the algebraic number $i$ does not change transcendence degree,
and
\[
 \Q(e,i\pi,i)=\Q(e,\pi,i).
\]
Therefore $\trdeg_{\Q}\Q(e,\pi)\ge2$, so $e$ and $\pi$ are
algebraically independent. The explicit adjoining of $i$ justifies the
transcendence-degree passage; the two smaller fields are not being
identified without explanation.

If $e+\pi$ were algebraic, there would be a nonzero $P\in\Q[T]$ with
$P(e+\pi)=0$. But $F(X,Y)=P(X+Y)$ is nonzero (set $Y=0$) and would satisfy
$F(e,\pi)=0$, a contradiction.\hfill$\square$

This proves
\[
 \boxed{\text{Schanuel's conjecture}\ \Longrightarrow\
        e+\pi\text{ is transcendental},}
\]
not an unconditional solution. The conjectural input remains part of the
hypothesis even though every subsequent step of the deduction is proved.

\paragraph{14. A more targeted route through $E$- and $G$-values.}
The following terminology is used in the sense of
\href{https://arxiv.org/html/2301.13518v1}{[E6]}, Section~1.
An $E$-function is a power series
\[
 f(z)=\sum_{n=0}^{\infty}\frac{a_n}{n!}z^n,\qquad a_n\in\Qbar,
\]
which satisfies a nonzero linear differential equation over $\Qbar(z)$,
and whose coefficient conjugates and common denominators grow at most
exponentially. More explicitly, for some $C>1$,
$|\sigma(a_n)|\le C^{n+1}$ for every embedding $\sigma$, and there exist
positive integers $d_n\le C^{n+1}$ such that $d_na_0,\ldots,d_na_n$ are
algebraic integers. A $G$-function has the same arithmetic conditions and
is a solution of a linear differential equation, but is written
$g(z)=\sum_{n\ge0}a_nz^n$ without factorials.

Let $\mathbf E$ be the ring of values of $E$-functions at algebraic points,
and let $\mathbf G$ be the ring of values at algebraic points of analytic
continuations of $G$-functions, using nonsingular determinations for the
values employed here. Both contain $\Qbar$. In particular,
\begin{equation}\label{eq:E-G-membership}
 e=\exp(1)\in\mathbf E,\qquad
 \pi=4\arctan(1)\in\mathbf G.
\end{equation}
The function $\arctan z$ is understood by continuation of its germ at zero
along the real segment to $1$. The point $1$ is not a singularity of this
germ's continuation; the nearest complex singularities are $i$ and $-i$.
Thus membership is not premised on absolute convergence of its Taylor
series at every point of the circle $|z|=1$.

If $s=e+\pi$ were algebraic, ring closure would imply
\[
 \pi=s-e\in\mathbf E,\qquad e=s-\pi\in\mathbf G,
\]
and consequently
\begin{equation}\label{eq:intersection-consequence}
 e,\pi\in\mathbf E\cap\mathbf G.
\end{equation}
Fischler and Rivoal explicitly discuss the conjecture
\begin{equation}\label{eq:intersection-conjecture}
 \boxed{\mathbf E\cap\mathbf G=\Qbar}
\end{equation}
and describe this property as out of reach in their discussion
\href{https://arxiv.org/html/2301.13518v1}{[E6]}.

\textbf{Proposition 5.} The intersection conjecture
\eqref{eq:intersection-conjecture} implies that $e+\pi$ is transcendental.

\emph{Proof.} If $e+\pi$ were algebraic, \eqref{eq:intersection-consequence}
and \eqref{eq:intersection-conjecture} would make $e$ and $\pi$ algebraic,
contrary to their known transcendence.\hfill$\square$

For this problem alone, either $e\notin\mathbf G$ or
$\pi\notin\mathbf E$ would already suffice. Calling this route
``more targeted'' refers to its formulation in terms of these particular
value rings; no implication or relative logical strength between the
full intersection conjecture and Schanuel's conjecture is asserted.

\paragraph{15. Mixed Hermite--Pad\'e proposal and its required certificate.}
The supplied development proposes rational polynomials
$\mathcal P_n,\mathcal Q_n,\mathcal T_n$ such that
\begin{equation}\label{eq:pade-function}
 F_n(z)=\mathcal P_n(z)+\mathcal Q_n(z)e^z+
          \mathcal T_n(z)\arctan z
\end{equation}
has high order of vanishing at zero. The calligraphic notation avoids
conflict with the convergents $P_j/Q_j$ of $\pi$. At $z=1$,
\begin{equation}\label{eq:pade-value}
 F_n(1)=\mathcal P_n(1)+\mathcal Q_n(1)e+
          \frac{\mathcal T_n(1)}4\pi.
\end{equation}
A finite homogeneous linear system can enforce prescribed initial Taylor
coefficients to vanish when there are enough polynomial coefficients.
That linear-algebra step is only a starting point. It does not by itself
control coefficient sizes, denominators, or $F_n(1)$.

\textbf{Correction to the supplied development.}
The supplied paragraph says that nonzero integer linear forms in
$1,e,\pi$ tending to zero would prove their rational linear independence.
That implication, without further rank or quantitative hypotheses, is
false. The already constructed forms
\[
 q_ne-p_n\ne0,\qquad |q_ne-p_n|\longrightarrow0
\]
are integer linear forms in $1,e,\pi$ with zero coefficient of $\pi$.
They would still exist if $\pi$ were replaced by $6-e$, although
$1,e,6-e$ are rationally dependent. Arbitrary small forms therefore do
not exclude the particular relation being tested.

Here is a sufficient certificate that does address the target.

\textbf{Proposition 6.} Suppose integers $U_n,V_n$ can be constructed for
infinitely many $n$ with
\begin{equation}\label{eq:aligned-certificate}
 \ell_n=U_n+V_n(e+\pi),\qquad
 \ell_n\ne0,\qquad |\ell_n|\longrightarrow0.
\end{equation}
Then $e+\pi$ is irrational.

\emph{Proof.} Under \eqref{eq:rationality},
$b\ell_n=bU_n+aV_n$ is a nonzero integer. Hence
$|\ell_n|\ge1/b$, contradicting convergence to zero.\hfill$\square$

To obtain this certificate from \eqref{eq:pade-value}, one sufficient
additional constraint is
\begin{equation}\label{eq:pade-alignment}
 \mathcal T_n(1)=4\mathcal Q_n(1).
\end{equation}
Choose a positive integer $D_n$ clearing the rational coefficients
occurring in \eqref{eq:pade-value}, and put
\[
 U_n=D_n\mathcal P_n(1),\qquad
 V_n=D_n\mathcal Q_n(1).
\]
Then $D_nF_n(1)=U_n+V_n(e+\pi)$. It remains to prove, for an infinite
sequence, both $D_nF_n(1)\ne0$ and $D_n|F_n(1)|\to0$.
Nonvanishing of the analytic function $F_n$ alone does not prove
nonvanishing of its value at $1$.

The missing ingredients are therefore coefficient-height control,
denominator control, analytic decay \emph{after} denominator clearing,
alignment such as \eqref{eq:pade-alignment}, and a nonvanishing argument
at the evaluation point. High-order vanishing at zero alone is especially
insufficient at $z=1$, which lies on the circle of convergence of the
Taylor expansion of $\arctan z$. A proof of full rational linear
independence of $1,e,\pi$ would need an appropriate independent family of
forms and a valid rank or determinant argument in addition.

This remains a research proposal, not a completed construction. The earlier
wording about the absence of a sufficiently strong mixed zero theorem is
read here as a description of the missing input in this attempt, not as a
proved nonexistence theorem or an exhaustive survey of all zero estimates.
Also the $E$-plus-$\arctan$ ansatz above should not be confused with every
technical use of ``mixed functions'' in \href{https://arxiv.org/html/2301.13518v1}{[E6]}.

\paragraph{16. Audit of the specific claimed proof cited in the development.}
Carella's version~2 of \emph{Linear Independence of Some Irrational
Numbers} is the manuscript cited in the supplied discussion
\href{https://arxiv.org/html/2007.15000v2}{[D1]}.
Its Section~7 assumes $A+Be+C\pi=0$ and invokes Lemma~5.2 to assert
$\sin(Be+A)\ne0$. Its Section~2 also assumes
$A-B\alpha=B\pi$ but treats $|\sin(A-B\alpha)|$ as positive.
These are the passages being examined, rather than an assertion that all
other work with a similar title has been checked.

Under the first assumed relation, one has exactly
\[
 Be+A=-C\pi,\qquad \sin(Be+A)=\sin(-C\pi)=0.
\]
For $t=Be+A$, the finite-sum bound involving $1/|\sin t|$ therefore cannot
be applied unless nonvanishing has been established independently. In
fact at $t\in\pi\Z$,
\[
 \frac1{2N}\sum_{n=-N}^{N}e^{2itn}
 =\frac{2N+1}{2N}\longrightarrow1\ne0.
\]

\textit{Logical qualification of the earlier criticism.}
A proved nonvanishing lemma contradicting a temporary hypothesis would be
a perfectly legitimate proof by contradiction. The issue is not merely
that the two statements conflict. The needed nonvanishing itself contains
the essential unresolved assertion. In the manuscript, Lemma~5.2 appeals
to Lemma~4.2. The displayed argument for Lemma~4.2 uses a comparison of the
form
\[
 x\ge y\quad\Longrightarrow\quad |x-c|\ge|y-c|,
\]
which is invalid in general: take $x=c=2$ and $y=1$. A lower bound on a
quantity does not remain a lower bound after subtracting a fixed number
and taking absolute values. In the actual displayed comparison, set
$k=m=1$ and $r=1$. Then $x=e-1<c=3\pi/2$, while for sufficiently large
convergent indices $j$,
$0<y=|q_je-p_j|<e-1$. Consequently $|x-c|<|y-c|$, the reverse of the
needed inequality. Here $p_j/q_j$ refers locally to the full sequence of
convergents used in that manuscript. Thus the cited chain does not provide
the independent nonvanishing input needed in Section~7.

The earlier Section~2 argument has a direct defect under its own
hypothesis:
\[
 A-B\alpha=B\pi\quad\Longrightarrow\quad
 |\sin(A-B\alpha)|=|\sin(B\pi)|=0.
\]
It cannot justify an eventual comparison with positive small sine errors
by declaring this fixed left-hand quantity positive. Positivity of
$A-B\alpha$ would not imply positivity of its sine's absolute value.
Consequently these arguments in the cited version do not establish the
claimed irrationality result. This is a diagnosis of the displayed
arguments, not a counterexample to the conjectured irrationality of
$e+\pi$.

\paragraph{17. Conversation context concerning formalization.}
The user's request to continue was motivated in part by Claude's reported
formalization of Fermat's Last Theorem. Anthropic's announcement dated
4 September 2026 describes a Lean formalization following the
Darmon--Diamond--Taylor exposition of Wiles's proof
\href{https://www.anthropic.com/research/formalizing-fermats-last-theorem}{[C1]}.
This context is retained, but neither the difficulty comparison nor the
success of a different formalization supplies a lemma about $e+\pi$.
The source statement is attributed to the announcement; this export does
not claim to have rechecked that Lean development. The earlier
continued-fraction response ended with a truncated final sentence; its
mathematical conclusion is fully recorded by Propositions~2--3 and the
explicitly unresolved condition \eqref{eq:missing-growth}.

\paragraph{18. Verification, corrections, and boundaries of the record.}
The denominator-matching proof includes an exact noncancellation argument.
The forbidden-interval proof uses the separation of distinct rational
numbers with explicitly bounded denominators. The denominator-counting
argument makes its finite initial exceptions and its dependence on $b$
explicit. The factorial calculations are identities and strict tail bounds,
not statistical observations about digits. The conditional transcendence
proofs retain their conjectural hypotheses.

The principal substantive correction is paragraph~15: small arbitrary
linear forms do not prove three-number linear independence. The additional
clarifications specify the factorial-expansion convention, combine the
thresholds $n\ge b$ and $n\ge3$, justify adjoining $i$ in the Schanuel
argument, distinguish nonvanishing of a function from nonvanishing at a
point, and explain why a contradiction with an independently proved lemma
would not itself be circular. The recurrence and digit converses are
spelled out as clarifications of the supplied reductions. No claim of a
new unconditional breakthrough is attached to these editorial additions.

As an export consistency check, the first thirty selected convergents were
independently generated from Euler's full continued-fraction digit pattern.
Exact integer arithmetic verified the initial fractions, the selected
recurrence, and the stated denominator bounds. The thirty error bounds
were also checked against a rational enclosure of $e$ obtained by
truncating its factorial series at $300$ and bounding the positive tail.
The finite telescoping identity behind factorial borrowing was checked
for $N=2,\ldots,20$. All these finite checks passed. They do not replace the
general written proofs above. In particular, such a check cannot establish $(*)$, the
infinitely-often growth bound \eqref{eq:missing-growth}, or an eventual
digit law for $\pi$. Numerical similarity between constants, or a finite
continued-fraction prefix, is not promoted to an equality or an
irrationality proof. The source list supplies background and dated status
statements; it does not supply any of the missing hypotheses identified
here.

\paragraph{19. Final assessment and exact remaining tasks.}
The principal unconditional implications established in this record are
\begin{gather}
 (*)\ \Longrightarrow\ e+\pi\notin\Q,\label{eq:summary-match}\\
 e+\pi\in\Q\ \Longrightarrow\
  \frac{\log Q_j}{j}\longrightarrow+\infty,\label{eq:summary-growth}\\
 e+\pi\in\Q\ \Longrightarrow\
  \{n!\pi\}=1-R_n\quad\text{for all sufficiently large }n.
 \label{eq:summary-factorial}
\end{gather}
The equivalent eventual factorial-digit and ceiling-recurrence
formulations have also been made explicit. In addition,
\[
 \text{Schanuel's conjecture}\ \Longrightarrow\ e+\pi\notin\Qbar,
 \qquad
 \mathbf E\cap\mathbf G=\Qbar\ \Longrightarrow\ e+\pi\notin\Qbar.
\]
No reverse implication between these conjectures is asserted.

An unconditional completion along the routes tried here would require a
proof of the matching condition $(*)$, or the infinitely-often exponential
denominator bound \eqref{eq:missing-growth}, or an obstruction to the exact
factorial fingerprint, or an appropriate special-case $E/G$ separation
result, or the aligned nonzero integer forms of
\eqref{eq:aligned-certificate}. None has been supplied for the actual
pair $(e,\pi)$.

The phrase ``mixed Lindemann--Weierstrass theorem'' in the supplied
conclusion is an informal description of a desired result controlling
\eqref{eq:mixed-exp-log}, not the name of an established theorem being
invoked. The existing identities control the separate pieces without
excluding the cancellation at issue.

\textbf{Outcome: exact reductions, conditional irrationality and
transcendence criteria, and identified proof obstructions;
UNRESOLVED.} This consolidated attempt establishes neither
$e+\pi\in\Q$ nor $e+\pi\notin\Q$ unconditionally.

\subsection{Sources}
The mathematical references below were checked on 27 September 2026.
References describing a conjecture, an unresolved status, or a claimed
proof are not endorsements of a solution. The user-supplied text and the
two earlier responses are the source of the research narrative;
the external references provide the specifically attributed background.

\begin{itemize}
 \raggedright
 \interlinepenalty=10000
 \item[S1] Piotr Kowalski, \emph{Positive characteristic Ax--Schanuel},
 arXiv:2311.00546v2, revised 11 March 2025, Section 2.1.
 \url{https://arxiv.org/html/2311.00546v2}.
 Original PDF link supplied by the user:
 \url{https://arxiv.org/pdf/2311.00546}.
 \textit{Explicit dated statement that irrationality of $e+\pi$ is unknown;
 also states Schanuel's conjecture.}

 \item[S2] Runlong Yu, \emph{Tail Criteria, No-Go Audits, and Ap\'ery-Type
 Certificate Obstructions for the Irrationality of $e+\pi$},
 arXiv:2606.17303, submitted 15 June 2026.
 \url{https://arxiv.org/abs/2606.17303}.
 \textit{Source retained from the first response. Abstract checked for its
 unresolved-status statement and description of factorial criteria.
 No independent validation of the preprint's entire audit is claimed.}

 \item[E1] Henry Cohn, \emph{A short proof of the simple continued fraction
 expansion of $e$}, American Mathematical Monthly 113 (2006), 57--62;
 arXiv:math/0601660v3.
 \url{https://arxiv.org/html/math/0601660v3}.
 \url{https://arxiv.org/abs/math/0601660}.
 \textit{Euler's continued fraction and the standard convergent recurrence.
 The selected-subsequence estimates and the conditional criteria are
 derived in the present record.}

 \item[E2] William Stein, \emph{Elementary Number Theory}, online course
 text, ``Convergence of Infinite Continued Fractions.''
 \url{https://wstein.org/edu/2007/spring/ent/ent-html/node65.html}.
 \textit{Continued-fraction approximation background; the error identity
 used here is also displayed explicitly in paragraph~2.}

 \item[E3] Boris Adamczewski and \'Eric Delaygue,
 \emph{Algebraic relations between sine and cosine values},
 arXiv:2406.01296v1, 3 June 2024, Theorem 3.
 \url{https://arxiv.org/html/2406.01296v1}.
 Original PDF link supplied by the user:
 \url{https://arxiv.org/pdf/2406.01296}.
 \textit{The algebraic-independence form of the Lindemann--Weierstrass
 theorem. This is applied only to algebraic exponents.}

 \item[E4] Vahagn Aslanyan, \emph{The Existential Closedness and Zilber--Pink
 Conjectures}, arXiv:2403.09304v1, 14 March 2024, Section 2.1.
 \url{https://arxiv.org/html/2403.09304v1}.
 \textit{Schanuel's conjecture and its application to the algebraic
 independence of $e$ and $\pi$ using $1$ and $i\pi$.}

 \item[E5] Michel Waldschmidt, \emph{Schanuel's Conjecture: algebraic
 independence of transcendental numbers}, Colloquium de Giorgi,
 6 May 2014, updated 20 May 2014, especially pages 1--4.
 \url{https://webusers.imj-prg.fr/~michel.waldschmidt/articles/pdf/ColloquiumDeGiorgiSchanuelConjecture2014.pdf}.
 \textit{Background on transcendence, Lindemann--Weierstrass, Schanuel's
 conjecture, and conditional consequences for familiar constants.}

 \item[E6] St\'ephane Fischler and Tanguy Rivoal,
 \emph{Relations between values of arithmetic Gevrey series, and
 applications to values of the Gamma function}, Journal of Number Theory
 261 (2024), 36--54; arXiv:2301.13518v1, Section 1.
 \url{https://arxiv.org/html/2301.13518v1}.
 Original PDF link supplied by the user:
 \url{https://arxiv.org/pdf/2301.13518}.
 \textit{Definitions of the function classes and value rings, and the
 discussion of $\mathbf E\cap\mathbf G=\Qbar$. The proposed
 $e^z$--$\arctan z$ construction is a proposal in the supplied development,
 not a theorem attributed to these authors.}

 \item[E7] Doron Zeilberger and Wadim Zudilin,
 \emph{The irrationality measure of $\pi$ is at most
 $7.103205334137\ldots$}, Moscow Journal of Combinatorics and Number Theory
 9 (2020), 407--419; arXiv:1912.06345v2.
 \url{https://arxiv.org/abs/1912.06345}.
 \textit{The stated irrationality-exponent bound. No exponential bound in
 the index of the convergents of $\pi$ is inferred from this result.}

 \item[D1] N. A. Carella, \emph{Linear Independence of Some Irrational
 Numbers}, arXiv:2007.15000v2, revision dated 27 February 2026.
 \url{https://arxiv.org/html/2007.15000v2}.
 \url{https://arxiv.org/abs/2007.15000v2}.
 \textit{Claimed proof examined, not an established result used as an
 input. The passages discussed are Section 2, Lemmas 4.2, 5.2 and 6.1,
 and Section 7. The version-specific arXiv record lists a later version;
 that later version is outside the scope of this export's audit.}

 \item[C1] Anthropic, \emph{Formalizing Fermat's Last Theorem},
 4 September 2026.
 \url{https://www.anthropic.com/research/formalizing-fermats-last-theorem}.
 \textit{Primary announcement retained only to document the conversation's
 motivating comparison. No theorem about $e+\pi$ is inferred from it.}

 \item[M1] User-provided LaTeX-format example,
 \texttt{Pasted text(20260927-133022).txt}, uploaded in this conversation.
 \textit{Formatting source: article preamble, unnumbered problem title,
 Definitions and mathematical statement, Short English statement,
 Research attempt, and Sources. Its BSD-specific metadata and mathematical
 assertions have not been reassigned to this problem.}
\end{itemize}
\end{document}
